\begin{definition}[The probabilists' Hermite polynomials]
For $k \in \mathbb{N}_0$ let $\mathrm{He}_k \in \mathbb{Z}[X]$ be the
probabilists' Hermite polynomials, defined by the recurrence
\[
  \mathrm{He}_0 = 1,
  \qquad
  \mathrm{He}_{k+1} = X\,\mathrm{He}_k - \mathrm{He}_k' ,
\]
where $p \mapsto p'$ is formal differentiation of polynomials.  Thus
$\mathrm{He}_1 = X$, $\mathrm{He}_2 = X^2 - 1$, and in general $\mathrm{He}_k$ is
monic of degree $k$ with integer coefficients.  This recurrence is the
definition of the family: it is the imported background object on which the
tablet's Hermite constructions are built, and nothing beyond the recurrence and
the two displayed initial values is assumed here.  These are the polynomials
orthogonal with respect to the weight $e^{-x^2/2}$ on $\mathbb{R}$; the
rescaling to the weight $e^{-x^2}$ used by the reference paper, and the
normalization of the resulting family, are carried out in the tablet's own
definition nodes and are not background.
\end{definition}
\begin{proposition}[Tonelli, Fubini, and Fubini for null sets]
Let $(\alpha,\mu)$ and $(\beta,\nu)$ be $\sigma$-finite measure spaces and let
$f : \alpha \times \beta \to \mathbb{R}$ be measurable for the product
$\sigma$-algebra.
\begin{enumerate}
\item (Tonelli) If $f \geq 0$ then
$\int f \, d(\mu\otimes\nu) = \int\!\!\int f(x,y)\, d\nu(y)\, d\mu(x)
 = \int\!\!\int f(x,y)\, d\mu(x)\, d\nu(y)$ as elements of $[0,+\infty]$.
\item (Fubini) If $f$ is $(\mu\otimes\nu)$-integrable then for $\mu$-a.e.\ $x$
the slice $y \mapsto f(x,y)$ is $\nu$-integrable, the function
$x \mapsto \int f(x,y)\,d\nu(y)$ is $\mu$-integrable, and the three integrals in
(1) agree as real numbers.
\item (Fubini for null sets) If $N \subseteq \alpha\times\beta$ is measurable
with $(\mu\otimes\nu)(N) = 0$, then for $\mu$-a.e.\ $x$ the slice
$N_x = \{y : (x,y) \in N\}$ satisfies $\nu(N_x) = 0$.
\end{enumerate}
\end{proposition}
\begin{proposition}[Dominated convergence]
Let $(\alpha,\mu)$ be a measure space, let $f_k : \alpha \to \mathbb{R}$ be
measurable with $f_k \to f$ $\mu$-a.e., and suppose there is a $\mu$-integrable
$g \geq 0$ with $|f_k| \leq g$ $\mu$-a.e.\ for every $k$.  Then $f$ and every
$f_k$ are $\mu$-integrable, $\int |f_k - f| \, d\mu \to 0$, and
$\int f_k \, d\mu \to \int f \, d\mu$.
\end{proposition}
\begin{proposition}[The scalar Gaussian integral]
$\int_{\mathbb{R}} e^{-x^2}\, dx = \sqrt{\pi}$, and more generally
$\int_{\mathbb{R}} e^{-b x^2}\, dx = \sqrt{\pi/b}$ for every real $b > 0$.
\end{proposition}
\begin{proposition}[Gaussian exponential moments]
For every $n \geq 1$ and every $t \in \mathbb{R}$,
\[
  \int_{\mathbb{R}^n} e^{t|x|}\, e^{-|x|^2}\, dx < \infty ,
\]
the integral being taken with respect to Lebesgue measure.  More generally
$\int_{\mathbb{R}^n} |P(x)|\, e^{t|x|}\, e^{-|x|^2}\, dx < \infty$ for every
polynomial $P$, since $|P(x)|e^{t|x|} \leq C_P e^{(t+1)|x|}$ for a constant
$C_P$.
\end{proposition}
\begin{proposition}[Cauchy--Schwarz]
Let $H$ be a real inner-product space.  Then
$|\langle x, y\rangle| \leq \|x\|\,\|y\|$ for all $x,y \in H$.  In integral
form: if $(\alpha,\mu)$ is a measure space and $f, g \in L^2(\mu)$ are real
valued, then the product $fg$ is $\mu$-integrable and
$\bigl|\int fg \, d\mu\bigr| \leq \|f\|_{L^2(\mu)}\,\|g\|_{L^2(\mu)}$.
\end{proposition}
\begin{proposition}[Jensen's inequality for a probability measure]
Let $\mu$ be a probability measure on $\alpha$, let $\varphi : \mathbb{R} \to
\mathbb{R}$ be convex, and let $f \in L^1(\mu)$ be real valued with
$\varphi \circ f \in L^1(\mu)$.  Then
$\varphi\bigl(\int f \, d\mu\bigr) \leq \int \varphi \circ f \, d\mu$.
\end{proposition}
\begin{proposition}[Riesz representation in a Hilbert space]
Let $H$ be a real Hilbert space and let $\Lambda : H \to \mathbb{R}$ be a
continuous linear functional.  Then there is a unique $y \in H$ with
$\Lambda(x) = \langle y, x\rangle$ for every $x \in H$, and
$\|\Lambda\|_{H^*} = \|y\|_H$.
\end{proposition}
\begin{proposition}[Uniqueness of characteristic functions]
Let $\mu$ and $\nu$ be finite Borel measures on $\mathbb{R}^n$ with
\[
  \int_{\mathbb{R}^n} e^{i\langle \xi, x\rangle}\, d\mu(x)
    = \int_{\mathbb{R}^n} e^{i\langle \xi, x\rangle}\, d\nu(x)
  \qquad\text{for every } \xi \in \mathbb{R}^n .
\]
Then $\mu = \nu$.
\end{proposition}
\begin{proposition}[Smoothness of a mollification, and the du Bois-Reymond lemma]
Let $f : \mathbb{R}^n \to \mathbb{R}$ be locally integrable for Lebesgue
measure.
\begin{enumerate}
\item If $\rho \in C_c^\infty(\mathbb{R}^n)$, then the convolution
\[
  (\rho * f)(x) = \int_{\mathbb{R}^n} \rho(t)\, f(x-t)\, dt
\]
is defined for every $x \in \mathbb{R}^n$ and belongs to
$C^\infty(\mathbb{R}^n)$.
\item If $\int_{\mathbb{R}^n} \varphi(x)\, f(x)\, dx = 0$ for every
$\varphi \in C_c^\infty(\mathbb{R}^n)$, then $f = 0$ Lebesgue-almost
everywhere.
\end{enumerate}
Item (2) is the fundamental lemma of the calculus of variations (du
Bois-Reymond).  No $L^p$ convergence of mollifications is asserted here; that is
a separate statement with a proof of its own, and it has its own tablet node.
\end{proposition}
\begin{proposition}[Borel--Cantelli]
Let $\mu$ be a measure on $\alpha$ and let $(A_k)_{k\in\mathbb{N}}$ be
measurable sets with $\sum_k \mu(A_k) < \infty$.  Then
$\mu\bigl(\limsup_k A_k\bigr) = 0$; equivalently, $\mu$-almost every point lies
in only finitely many $A_k$.
\end{proposition}
\begin{proposition}[Fatou's lemma]
Let $(\alpha,\mu)$ be a measure space and let $f_k : \alpha \to [0,+\infty]$ be
measurable.  Then
$\int \liminf_k f_k \, d\mu \leq \liminf_k \int f_k \, d\mu$.
\end{proposition}
\begin{proposition}[Entire functions and uniqueness of power-series coefficients]
Let $f : \mathbb{C} \to \mathbb{C}$ be entire.  Then $f$ is represented on all
of $\mathbb{C}$ by its Taylor series at $0$,
$f(z) = \sum_{k\geq0} c_k z^k$ with $c_k = f^{(k)}(0)/k!$, the series
converging absolutely and locally uniformly.  Moreover the coefficients are
unique: if $\sum_k c_k z^k = \sum_k d_k z^k$ for all $z$ in some neighbourhood
of $0$, with both series convergent there, then $c_k = d_k$ for every $k$.
\end{proposition}
\begin{proposition}[$\sigma$-finiteness of Lebesgue measure]
Lebesgue measure on $\mathbb{R}^n$ is $\sigma$-finite: $\mathbb{R}^n$ is the
union of the countably many balls $B(0,k)$, $k \in \mathbb{N}$, each of finite
Lebesgue measure.  The same therefore holds for any measure obtained from
Lebesgue measure by multiplying by a locally bounded density.
\end{proposition}
\begin{proposition}[Density of compactly supported continuous functions in $L^p$]
Let $\mu$ be a regular Borel measure on $\mathbb{R}^n$ (Lebesgue measure is
one), let $1 \leq p < \infty$ and let $f \in L^p(\mu)$ be real valued.  For
every $\eta > 0$ there is $g : \mathbb{R}^n \to \mathbb{R}$ which is
continuous, has compact support, lies in $L^p(\mu)$, and satisfies
$\|f - g\|_{L^p(\mu)} \leq \eta$.  Any such $g$ is bounded and uniformly
continuous on $\mathbb{R}^n$: it is continuous, and if $S$ is a compact set
containing $\operatorname{supp} g$ in its interior then $g$ is uniformly
continuous on the compact set $S$ by the Heine--Cantor theorem and vanishes
identically outside $\operatorname{supp} g$, which gives a single modulus of
continuity valid on all of $\mathbb{R}^n$.
\end{proposition}
\begin{proposition}[Translation invariance of Lebesgue measure]
Let $\lambda$ denote Lebesgue measure on $\mathbb{R}^n$ and let
$t \in \mathbb{R}^n$.  Then $\lambda(A + t) = \lambda(A)$ for every Lebesgue
measurable $A \subseteq \mathbb{R}^n$, and consequently, for every measurable
$h : \mathbb{R}^n \to [0,+\infty]$ and every measurable $A$,
\[
  \int_A h(x-t)\, dx = \int_{A-t} h(y)\, dy ,
\]
both sides being equal in $[0,+\infty]$; the same identity holds for
$\lambda$-integrable real-valued $h$.  Lebesgue measure is also invariant under
the reflection $x \mapsto -x$, so $\int h(-x)\,dx = \int h(x)\,dx$ for every
such $h$; combining the two invariances, the map $t \mapsto x - t$ preserves
$\lambda$ for every fixed $x$, and $\int h(x-t)\,d\lambda(t) = \int h\,d\lambda$.
\end{proposition}
\begin{proposition}[Summable families of reals, and the extended-real bridge]
Let $I$ be a countable index set and let $(c_i)_{i\in I}$, $(d_i)_{i \in I}$ be
families of real numbers.
\begin{enumerate}
\item (Comparison) If $0 \leq c_i \leq d_i$ for every $i$ and $(d_i)$ is
summable, then $(c_i)$ is summable and $\sum_i c_i \leq \sum_i d_i$.
\item (Linearity) If $(c_i)$ and $(d_i)$ are summable and $a, b \in \mathbb{R}$,
then $(a c_i + b d_i)_{i}$ is summable with
$\sum_i (a c_i + b d_i) = a\sum_i c_i + b \sum_i d_i$.
\item (Extended-real bridge) Suppose $c_i \geq 0$ for every $i$, and read
$\sum_i c_i$ as an element of $[0,+\infty]$ by taking the supremum of the finite
partial sums.  Then $\sum_i c_i < \infty$ if and only if $(c_i)$ is summable as a
family of real numbers, and in that case the two readings of $\sum_i c_i$ agree.
\item (Vanishing) If $e_i \in [0,+\infty]$ for every $i$ and
$\sum_i e_i = 0$, then $e_i = 0$ for every $i$; this is immediate from
$e_{i_0} \leq \sum_i e_i$ for each fixed $i_0$, the terms being nonnegative.
\end{enumerate}
Clause (1) applied together with the pointwise inequality
$(s+t)^2 \leq 2s^2 + 2t^2$ and clause (2) shows that, for any fixed weights
$w_i \geq 0$, the families $(c_i)$ with $\sum_i w_i c_i^2 < \infty$ form a linear
subspace of $\mathbb{R}^I$.
\end{proposition}
\begin{proposition}[Weak-$*$ sequential compactness of dual balls, separable case]
Let $E$ be a normed space over $\mathbb{R}$ which is \emph{separable}, i.e.\ has
a countable dense subset, and let $E^*$ denote its dual: the continuous linear
functionals $E \to \mathbb{R}$, normed by
$\|\Lambda\|_{E^*} = \sup_{\|w\| \leq 1} |\Lambda(w)|$.  Let $R \geq 0$ and let
$(\Lambda_k)_{k \geq 1}$ be a sequence in $E^*$ with $\|\Lambda_k\|_{E^*} \leq R$
for every $k$.  Then there are indices $k_1 < k_2 < \cdots$ and
$\Lambda \in E^*$ with $\|\Lambda\|_{E^*} \leq R$ such that
\[
  \Lambda_{k_j}(w) \longrightarrow \Lambda(w)
  \qquad\text{as } j \to \infty, \text{ for every } w \in E .
\]
The two hypotheses are exactly separability of $E$ and completeness of the
scalar field; nothing is asserted about sequential compactness in $E$ itself, and
no reflexivity is used.  Transporting this statement to a Hilbert space, and
verifying separability of the Hilbert space in question, are not part of this
item.
\end{proposition}
\begin{proposition}[Elementary facts about $\liminf$]
Let $(x_k)$, $(y_k)$, $(z_k)$ be sequences.
\begin{enumerate}
\item (Monotonicity) If the terms take values in $[0,+\infty]$, or if they are
real and both sequences are bounded, and $x_k \leq y_k$ for every $k$, then
$\liminf_k x_k \leq \liminf_k y_k$.
\item (Convergent sequences) If $x_k \to x$, with values in $\mathbb{R}$ or in
$[0,+\infty]$, then $\liminf_k x_k = x$.
\item (Adding a convergent sequence) If $(y_k)$ is a bounded real sequence and
$z_k \to z$ in $\mathbb{R}$, then $(y_k + z_k)$ is bounded and
$\liminf_k (y_k + z_k) = \bigl(\liminf_k y_k\bigr) + z$.
\end{enumerate}
\end{proposition}
\begin{proposition}[Monotonicity and subadditivity of a measure]
Let $\mu$ be a measure on a measurable space $(\alpha, \mathcal{A})$.
\begin{enumerate}
\item (Monotonicity) If $A \subseteq B$ then $\mu(A) \leq \mu(B)$.
\item (Subadditivity) $\mu(A \cup B) \leq \mu(A) + \mu(B)$ for all $A, B$, and
more generally $\mu\bigl(\bigcup_{k} A_k\bigr) \leq \sum_k \mu(A_k)$ for any
countable family $(A_k)$.
\item (Monotonicity up to null sets) If $A \subseteq B \cup N$ with
$\mu(N) = 0$, then $\mu(A) \leq \mu(B)$; this follows from (1) and (2), since
$\mu(A) \leq \mu(B \cup N) \leq \mu(B) + \mu(N) = \mu(B)$.
\end{enumerate}
All assertions are as elements of $[0,+\infty]$.
\end{proposition}
\begin{proposition}[Compact sets: extreme values and finite Lebesgue measure]
\begin{enumerate}
\item (Extreme value theorem) Let $K$ be a nonempty compact subset of
$\mathbb{R}^n$ and let $\varrho : \mathbb{R}^n \to \mathbb{R}$ be continuous.
Then $\varrho$ attains both a minimum and a maximum on $K$: there are
$x_-, x_+ \in K$ with $\varrho(x_-) \leq \varrho(x) \leq \varrho(x_+)$ for every
$x \in K$.
\item (Finite Lebesgue measure) Every compact $K \subseteq \mathbb{R}^n$ is
bounded, hence $K \subseteq B(0,r)$ for some $r \in \mathbb{N}$, and therefore
$\lambda(K) \leq \lambda(B(0,r)) < \infty$ by monotonicity of $\lambda$ and the
finiteness of the measure of a ball.
\end{enumerate}
\end{proposition}
